% !TEX root = ../main.tex
\hspace{24pt}

\section{Motivation}

Traditional Chinese Painting (TCP) is a classical form of Chinese visual art,
of which ink wash painting is the most renowned;
it conveys distance and a sense of space through ink-tone variation and the deliberate use of white space.
A painting, however, is inherently a static two-dimensional image and lacks any sense of motion.
In today's digital media environment, audience expectations have shifted from static toward dynamic content,
and social platforms as well as digital exhibitions increasingly present artworks in video form to reach a broader audience.
Presenting TCP as animated videos would not only allow viewers to perceive the spatial layering of a painting more intuitively,
but would also help traditional art reach younger generations in the digital era and extend its cultural value.
In recent years, numerous studies have addressed 3D model reconstruction of real-world objects \cite{zero123}\cite{difix3d}
and scene-level 2.5D video presentation \cite{viewcrafter}\cite{wonderworld}\cite{trajectorycrafter}.
A 2.5D video is a visual representation that lies between a flat 2D image and a full 3D reconstruction:
driven by the depth information of a single image, a virtual camera moves along a predefined trajectory,
producing an animated video with parallax that conveys spatial depth without constructing a complete geometric model.
Since digitized TCP data are far scarcer than real-world imagery,
full 3D reconstruction---which requires large multi-view datasets or precise geometric annotations---faces
prohibitively high barriers in the TCP domain.
By contrast, diffusion-model-based techniques can synthesize a 2.5D video from a single image,
making convincing visual results much more attainable under limited data.
Applying such techniques to TCP would allow traditional art to be presented in a more engaging and immersive manner.
Prior work such as ViewCrafter \cite{viewcrafter} has explored synthesizing parallax videos from a single image with diffusion models
and demonstrated compelling results on real-world photographs, confirming the feasibility of this direction.
However, ViewCrafter is trained primarily on real photographs;
when confronted with TCP-specific characteristics---white-space composition, abstract brushstrokes,
and the absence of photometric depth cues---it tends to produce unnatural depth estimates and image distortions,
and therefore cannot be applied directly.

Although modern diffusion models possess strong visual semantic priors
acquired from large-scale and wide-coverage training data \cite{LAION-5B},
generating a 2.5D video directly from a single TCP image remains highly problematic.
In the prior knowledge of these models, TCP material appears almost exclusively in the form of flat images,
so a model tends to interpret the input as a flat object hanging on a wall;
the generated video then merely shows a camera orbiting around a picture frame
rather than moving through the three-dimensional space of the scene.
Accurate depth information must therefore be obtained before video generation,
so that the model can understand the true spatial structure of the painting.

Meanwhile, prevailing diffusion-based monocular depth estimation methods are trained on
large collections of real photographs and are prone to notable misjudgments
when handling the white space and abstract brushstrokes unique to TCP.
The visual language of TCP differs fundamentally from that of photographs:
large blank regions in a painting represent distant scenery,
yet a model may interpret them as nearby empty planes according to its geometric priors;
regions of dense ink do not necessarily indicate nearby objects,
but rather the painter's expressive emphasis.
Owing to these differences, directly applying existing models such as Marigold \cite{marigold}
often yields depth maps inconsistent with the actual spatial relations of the painting.
Furthermore, datasets of authentic TCP with depth annotations are difficult to obtain,
and conventional manual annotation cannot accumulate a sufficient amount of data,
rendering both fine-tuning existing models and training a TCP-specific depth estimation
diffusion model from scratch infeasible.

\section{Research Objectives and Contributions}

To address the above problems, this study adopts Marigold-DC \cite{marigolddc},
which allows users to guide the estimation with a small number of sparse depth control points (sparse depth hints)
without dense full-image annotation, thereby greatly reducing manual effort.
In addition, the loss function of the original Marigold-DC measures geometric accuracy only.
Taking into account the brushstroke structure and white-space semantics unique to TCP,
this study improves the loss function and designs \textit{ChinesePaintingLoss},
a composite loss tailored to the characteristics of Chinese painting (detailed in Section~\ref{sec:depth}),
so that the estimated depth better follows the spatial conventions of TCP.

\begin{figure}[htbp]
  \centering
  \includegraphics[width=\textwidth]{pipeline}
  \caption{The proposed 2.5D video generation pipeline for Traditional Chinese Painting.
  Given a single TCP input image, depth estimation first yields a rough depth map $D_{\mathrm{rough}}$;
  based on $D_{\mathrm{rough}}$, the user annotates sparse depth control points $G$ (manual adjustment),
  which serve as guidance for the improved Marigold-DC to output a refined depth map $D_{\mathrm{refined}}$;
  $D_{\mathrm{refined}}$ is then unprojected into a three-dimensional point cloud ($\mathit{PCD}$) via Depth Anything V3;
  finally, Uni3C renders the point-cloud video $V_{\mathit{pcd}}$ along a predefined camera trajectory
  and drives a video diffusion model to generate a smooth 2.5D video.}
  \label{fig:pipeline}
\end{figure}

Although the injected sparse depth control points yield a relatively accurate depth map,
directly feeding the depth map together with a camera trajectory into a video diffusion model
still suffers from insufficient geometric information.
A depth map is two-dimensional data expressed in image space;
it cannot be used directly to render the scene from novel viewpoints,
and the model lacks sufficient three-dimensional geometric cues
to reason about occlusion and displacement of objects across viewpoints,
so the generated video lacks convincing parallax.

Therefore, after obtaining accurate TCP depth values through the proposed estimation pipeline,
we unproject the depth map into a three-dimensional point cloud colored by the original image.
The point cloud explicitly records the geometric position of every point in 3D space
and can be rendered along a predefined camera trajectory using standard rasterization,
producing a preliminary 2.5D point-cloud video that is rich in geometric information and exhibits true parallax.
Although the rendered result contains occlusion holes exposed by viewpoint shifts,
it provides the subsequent video diffusion model with explicit spatial-geometric conditions.
Finally, this study designs text prompts dedicated to the TCP domain
to steer the diffusion model toward completions consistent with the ink-and-brush style,
and automatically computes the parameters of a circular camera trajectory
from the painting size and the point-cloud geometry,
ensuring that the magnitude and speed of viewpoint shifts match the visual rhythm of TCP.
These conditions, together with the point-cloud video, are fed into Uni3C \cite{uni3c},
whose video diffusion model fills the holes under geometric guidance
and outputs a smooth and complete 2.5D video.

The goal of this study is to establish TCP-Uni3C---a complete end-to-end pipeline
with which a user provides only a single TCP image and a small number of manual depth annotations
to obtain a 2.5D video with a convincing sense of depth,
enabling traditional Chinese paintings to be presented in animated form.
The overall pipeline is illustrated in Figure~\ref{fig:pipeline}.

It is worth noting that this pipeline is not a purely automatic black-box system;
instead, it provides user controls at multiple levels:
users may intervene in the spatial semantics via sparse depth control points,
design the camera motion via trajectory parameters (e.g., rotation angles and translation distances),
and guide the video diffusion model toward the desired visual style via text prompts,
giving the final video a considerable degree of controllability.

The main contributions of this study are summarized as follows:

\begin{itemize}
  \item To correct the estimation bias of existing diffusion-based depth estimation models on TCP-style images,
        we design an interactive sparse-control-point annotation interface
        with which users click positions directly on the image and assign depth values;
        the annotations are conveniently fed into Marigold-DC as guidance conditions,
        so that a small number of manual hints effectively rectify erroneous regions in the depth map.

  \item We design \textit{ChinesePaintingLoss}, a composite loss function tailored to the characteristics of TCP,
        which augments the geometric-accuracy constraint with a brushstroke-aware smoothness loss
        and a whitespace regularization term,
        so that the refined depth map simultaneously satisfies the numerical constraints of the control points
        and the spatial-semantic conventions of Chinese painting.

  \item We propose TCP-Uni3C---a complete end-to-end pipeline that generates a 2.5D video
        from a single TCP image and a small number of manual depth annotations,
        without constructing 3D models or collecting large-scale TCP training datasets.
\end{itemize}
